\chapter{Basic quantum mechanics and electrodynamics}
\label{cha:appendix_basic}

In this appendix, we develop the basic quantum-mechanical and electromagnetic descriptions used in chapters~\ref{cha:fundamentals_a} and~\ref{cha:fundamentals_b}.
We first discuss quantum time evolution and the expectation values of many-body observables.
We then examine the Thomas--Fermi construction and its extensions.
Finally, we develop the electromagnetic constitutive relations and the role of causality in the optical response.
The working density-functional and optical relations remain in the corresponding chapters.

The quantum-mechanical and density-functional discussion partly draws on my notes from the Coursera course \textit{Density Functional Theory}, offered by École Polytechnique and taught by Francesco Sottile and Lucia Reining~\cite{sottile_reining_dft_coursera}.
We retain the Hartree atomic units introduced in section~\ref{sec:manybody_schrodinger_equation}, unless constants are restored explicitly.
The electrodynamic discussion uses SI units, as in chapter~\ref{cha:fundamentals_b}.

\newpage
\section{A story starts with time evolution quantum mechanics}
\label{sec:time_evolution_quantum_mechanics}

% opening paragraph
% Imagine that we are watching a microscopic play, with electrons and atoms as the actors.
Imagine that we are watching the evolution of electrons and atoms as a microscopic performance.
You blink, and suddenly the stage configuration has shifted.
Something has changed. But how do we describe this evolution in physical language?

In classical physics, we started with Newton's laws to describe motion.
Later, Lagrangian and Hamiltonian mechanics provided a different framework for describing physics from the perspective of degrees of freedom, symmetries, and conservation laws.
The twentieth century then brought the quantum revolution, establishing quantum mechanics as the fundamental framework for describing microscopic physics.
In this new framework, the Schrödinger equation plays the central role.

\subsection{Time-dependent Schrödinger equation}

In quantum mechanics, the state of a physical system at a given instant is represented by a vector $\ket{\psi(t)}$.
In a chosen representation, the same state can be expressed as a wave function.
This state vector encapsulates all information required to predict the statistics of measurement outcomes at time $t$.

Next, we switch to the coordinate representation.
We define the wave function as the projection of the state onto the position eigenbasis:

\begin{equation}
    \psi(\vec{r}\,,t)\coloneqq \braket{\vec{r}\,|\psi(t)}
\label{coordinate_representation_definition}
\end{equation}

here, $\ket{\vec{r}\,}$ is an eigenstate of the position operator $\hat{\vec{r}}$.

The time evolution of the state is determined by the hamiltonian operator $\hat{H}$.
It corresponds to the total energy of the system.
For a closed system with a time-independent hamiltonian, the expectation value of the energy is conserved. This evolution is described by the time-dependent Schrödinger equation~\cite{shankar2012principles}:

\begin{equation}
    i\hbar\frac{\partial}{\partial t}\ket{\psi(t)}=\hat{H}\ket{\psi(t)}
\label{time-dependent_Schrodinger_equation}
\end{equation}

here, the left-hand side gives the time derivative of the state vector;
the right-hand side shows that this change is generated by the action of the hamiltonian operator.

Suppose we know the state vector at an initial time $t=t_0$. This elegant equation then uniquely determines its future state and, in turn, the probabilities of measurement outcomes.

\subsection{Time evolution operator}

Suppose that the hamiltonian is independent of time, so that it does not depend explicitly on $t$.
In this case, the Schrödinger equation admits a formal solution.
Given a reference state vector $\ket{\psi(t_0)}$,
the state at time $t$ can be expressed with the time-evolution operator $\hat{U}(t,t_0)$:

\begin{equation}
    \ket{\psi(t)}=\hat{U}(t,t_0)\ket{\psi(t_0)}
\label{state_time_evolution}
\end{equation}

Subsequently, we introduce the explicit exponential representation of the time-evolution operator:

\begin{equation}
    \hat{U}(t,t_0)=\mathrm{e}^{-\frac{\mathrm{i}}{\hbar}\hat{H}(t-t_0)}
\label{time_evolution_operator}
\end{equation}

With the time difference $\tau=t-t_0$, we can express the time-evolution as a single-parameter family $\hat{U}(\tau)$:

\begin{equation}
    \hat{U}(\tau)=\mathrm{e}^{-\frac{\mathrm{i}}{\hbar}\hat{H}\tau}
\label{time_evolution_operator_tau}
\end{equation}

If we stipulate $t_0=0$, this expression becomes:

\begin{equation}
    \hat{U}(t)=\mathrm{e}^{-\frac{\mathrm{i}}{\hbar}\hat{H}t}
\label{time_evolution_operator_reduced_form}
\end{equation}

We use $\hat{U}(t,t_0)$ to propagate the state forward in time while preserving its norm.
Since $\hat{H}$ is Hermitian, $\hat{U}(t,t_0)$ is unitary.
This property can be described as:

\begin{equation}
    \hat{U}^\dagger(t,t_0)\,\hat{U}(t,t_0)=\hat{U}(t,t_0)\,\hat{U}^\dagger(t,t_0) = \hat{I}
\label{evolution_unitarity}
\end{equation}

Starting from the unitarity of $\hat{U}(t,t_0)$, quantum time evolution is reversible.
Here, we denote the evolution from time $t$ back to the reference time $t_0$ as $\hat{U}(t_0,t)$:

\begin{equation}
    \hat{U}^{\dagger}(t,t_0)= \hat{U}^{-1}(t,t_0) = \hat{U}(t_0,t)
\label{time_evolution_reversible}
\end{equation}

In the present time-independent case, we can further show that the Hermitian conjugate corresponds to evolving backward in time.
In terms of the time difference $\tau$, this relation can be written as:

\begin{equation}
    \hat{U}^{\dagger}(\tau)= \hat{U}^{-1}(\tau) = \hat{U}(-\tau)
\label{time_evolution_reversible_tau}
\end{equation}

% Accordingly, in the stipulation of $t=t_0$, this expression degenerates into:
Accordingly, if we choose $t_0=0$, this expression becomes:

\begin{equation}
    \hat{U}^{\dagger}(t)= \hat{U}^{-1}(t)= \hat{U}(-t)
\label{time_evolution_reversible_deg}
\end{equation}

The two-time evolution operator $\hat{U}(t,t_0)$ maps a quantum state from the reference time $t_0$ to time $t$.
By construction, successive time evolutions must be consistent with each other.
This requirement leads to the composition law:

\begin{equation}
    \hat{U}(t_2,t_0)=\hat{U}(t_2,t_1)\hat{U}(t_1,t_0)
\label{composition_law_time_evolution}
\end{equation}

In the time-independent case, $\hat{U}(t,t_0)$ depends only on the time difference $\tau=t-t_0$.
The operators $\hat{U}(\tau)$ therefore form a one-parameter unitary group:

\begin{equation}
    \hat{U}(\tau_2)\,\hat{U}(\tau_1)=\hat{U}(\tau_2+\tau_1)
\label{time_evolution_group_property_tau}
\end{equation}

In particular, when we choose $t_0=0$, the one-parameter composition law becomes:

\begin{equation}
    \hat{U}(t_1+t_2)=\hat{U}(t_1)\,\hat{U}(t_2)
\label{time_evolution_group_property}
\end{equation}

For the one-parameter family $\hat{U}(\tau)$, the identity element, $\hat{I}$, corresponds to $\tau=0$:

\begin{equation}
    \hat{U}(0)=\hat{I}
\label{time_evolution_identity}
\end{equation}

In this sense, these relations mean that time translation can be composed consistently, and that the evolution can always be undone by evolving backward.

Unitarity also guarantees probability conservation:

\begin{equation}
\begin{aligned}
    \braket{\psi(t)|\psi(t)}
    &= \bra{\psi(t_0)}\hat{U}^{\dagger}(t,t_0)\,\hat{U}(t,t_0)\ket{\psi(t_0)} \\
    &= \bra{\psi(t_0)}\hat{I}\ket{\psi(t_0)} \\
    &= \braket{\psi(t_0)|\psi(t_0)}
\end{aligned}
\label{time_evolution_norm_conservation}
\end{equation}

If the state is normalized at time $t_0$, then $\braket{\psi(t_0)|\psi(t_0)}=1$, and the state remains normalized for all later times.

This is why we can view $\hat{U}(t,t_0)$ as the quantum analogue of a time-translation operator.
It moves the state through Hilbert space without changing its total probability.

Finally, for a small time step $\Delta t$, we return to the exponential form:

\begin{equation}
    \hat{U}(\Delta t)=\mathrm{e}^{-\frac{\mathrm{i}}{\hbar}\hat{H}\Delta t}
\label{time_evolution_short_time_exact}
\end{equation}

and obtain the short-time expansion:

\begin{equation}
    \hat{U}(\Delta t)\approx \hat{I}-\frac{\mathrm{i}}{\hbar}\hat{H}\Delta t+\mathcal{O}(\Delta t^2)
\label{time_evolution_short_time_series}
\end{equation}

Acting on $\ket{\psi(t)}$ and taking the limit $\Delta t\to 0$ then recovers the differential form of the time-dependent Schrödinger equation.

% \subsection{Generalizing to time-dependent hamiltonians}

We now consider the more general case in which the hamiltonian depends explicitly on time and is denoted by $\hat{H}(t)$.
The state is still propagated by the evolution operator $\hat{U}(t_1,t_0)$ as in equation~\eqref{state_time_evolution}.

In this case, its corresponding time-evolution operator can be written as a time-ordered exponential~\cite{sakurai2020modern}:

\begin{equation}
    \hat{U}(t_1,t_0)=\hat{T}\exp\left(-\frac{\mathrm{i}}{\hbar}\int_{t_0}^{t_1}\hat{H}(t)\mathrm{d}t\right)
\label{time-ordering_operator}
\end{equation}

here, $\hat{T}$ denotes the time-ordering operator.
It orders products of $\hat{H}(t)$ such that operators at later times appear to the left.
This is required when the hamiltonians at different times do not commute:

\begin{equation}
    [\hat{H}(t),\hat{H}(t')]\neq 0
\label{hamiltonian_noncommutativity}
\end{equation}

Equivalently, the hamiltonian operators evaluated at earlier times act first on the state vector.

The corresponding many-body hamiltonian is introduced in section~\ref{sec:manybody_schrodinger_equation}.
We now return to the time evolution and its representation.

\subsection{From Schrödinger picture to Heisenberg picture}

Up to now, we have adopted the Schrödinger picture to describe the time evolution of quantum systems. In this picture, states carry all the time dependence, while operators remain fixed. Specifically, the quantum state $\ket{\psi(t)}$ evolves in time according to the Schrödinger equation~\eqref{time-dependent_Schrodinger_equation}.

For a given observable represented by a Hermitian (self-adjoint) operator $\hat{\Omega}$, its expectation value is calculated as:

\begin{equation}
    \braket{\hat{\Omega}} = \braket{\psi(t)|\hat{\Omega}|\psi(t)}
\label{expectation_Hermitian}
\end{equation}

In the standard formulation of quantum mechanics, observables represent measurable physical quantities, so their expectation values are real.
This is why observables are taken to be Hermitian operators, which ensures real eigenvalues and guarantees that $\braket{\hat{\Omega}}\in\mathbb{R}$.

To set up the transformation, we return to the Schrödinger picture.
Recall \eqref{state_time_evolution} and the definition of the time-evolution operator in \eqref{time_evolution_operator}.
With $\hat{U}(t,t_0)$, the two pictures are related as follows.

In the Heisenberg picture, the state at the reference time $t_0$ is taken to be time-independent:

\begin{equation}
    \ket{\psi_\mathrm{H}}=\ket{\psi(t_0)}
\label{heisenberg_state_definition}
\end{equation}

We keep the state fixed and shift the time dependence to the operators.
For any Schrödinger picture operator $\hat{\Omega}$, we can define its Heisenberg counterpart $\hat{\Omega}_\mathrm{H}(t)$ as:

\begin{equation}
    \hat\Omega_\mathrm{H}(t)
    =\hat{U}^\dagger(t,t_0)\,\hat{\Omega}\,\hat{U}(t,t_0)
    =\mathrm{e}^{\frac{\mathrm{i}}{\hbar}\hat{H}(t-t_0)}\,\hat{\Omega}\,\mathrm{e}^{-\frac{\mathrm{i}}{\hbar}\hat{H}(t-t_0)}
\label{heisenberg_operator_definition}
\end{equation}

This definition transfers the time dependence from the state to the operator.

With these definitions, the expectation value is unchanged from Schrödinger to the Heisenberg picture.
Applying the expression of time-evolution operator\eqref{state_time_evolution},
the following relation demonstrates the equivalence of the two pictures:

\begin{equation}
\begin{aligned}
    \braket{\hat{\Omega}}(t)&=\bra{\psi(t)}\hat{\Omega}\ket{\psi(t)}\\
    &=\bra{\psi(t_0)}\hat{U}^\dagger(t,t_0)\,\hat{\Omega}\,\hat{U}(t,t_0)\ket{\psi(t_0)} \\
    &=\bra{\psi_\mathrm{H}}\hat{\Omega}_\mathrm{H}(t)\ket{\psi_\mathrm{H}}
\end{aligned}
\label{equivalence_schrodinger_heisenberg}
\end{equation}

For a time-independent hamiltonian and an operator without explicit time dependence in the Schrödinger picture, differentiating \eqref{heisenberg_operator_definition} gives the Heisenberg equation of motion~\cite{shankar2012principles}:

\begin{equation}
    \frac{\mathrm{d}}{\mathrm{d}t}\hat{\Omega}_\mathrm{H}(t)
    =\frac{\mathrm{i}}{\hbar}\left[\hat{H},\hat{\Omega}_\mathrm{H}(t)\right]
\label{heisenberg_equation_of_motion}
\end{equation}

This form is particularly convenient for discussing conserved quantities and symmetries, since it expresses the dynamics directly at the operator level.

\subsection{From Heisenberg picture to the continuity equation}

Having seen that any local one-body observable can be written as an integral over a single function,
we now introduce that central quantity.
The one-body density contains all the information needed to compute local one-body expectation values and serves as the fundamental variable in density functional theory~\cite{hohenberg1964inhomogeneous,jones2015density,maitra2016perspective}.
In what follows, we review its precise definition:

\begin{equation}
    n(\vec{r}\,)
    = \bra{\psi}\,\hat{n}(\vec{r}\,)\ket{\psi}
    = \bra{\psi}\sum_i\delta(\vec{r} - \hat{\vec{r}}_i\,)\ket{\psi}
\label{particle_density_operator_definition}
\end{equation}

This defines the electron density at point $\vec{r}$.
Physically, $n(\vec{r}\,)\mathrm{d}^3r$ is the expected number of electrons in the volume element $\mathrm{d}^3r$ around $\vec{r}$.
In density functional theory, $n(\vec{r}\,)$ serves as the fundamental variable replacing the full many-body wave function.  

It can be shown that the ground-state density $n(\vec{r}\,)$ uniquely determines the external potential up to an additive constant, and hence all ground-state observables can be written as functionals of $n(\vec{r}\,)$.
This dramatic reduction comes from a 3$N$-dimensional wave function to a three-dimensional density, which lies at the heart of density functional theory, making formerly intractable many-body problems computationally feasible.

The time-dependent density-functional theory relies on the continuity equation in the Heisenberg picture.
We first define the particle density operator in the Schrödinger-picture as:

\begin{equation}
    \hat{n}(\vec{r}\,)
    = \sum_i \delta(\vec{r}-\hat{\vec{r}}_i\,)
\label{Schrodinger-picture_density_operator}
\end{equation}

and its corresponding Heisenberg-picture counterpart is defined by the unitary time evolution as:

\begin{equation}
    \hat{n}_\mathrm{H}(\vec{r},t)
    = \hat U^\dagger(t)\,\hat{n}(\vec{r}\,)\,\hat{U}(t)
\label{Heisenberg-picture_density_operator}
\end{equation}

Its time derivative follows from the Heisenberg equation of motion,

\begin{equation}
    \frac{\partial}{\partial t}\hat n_\mathrm{H}(\vec{r},t)
    = \frac{\mathrm{i}}{\hbar}\left[\hat{H},\hat n_\mathrm{H}(\vec{r},t)\right]
\label{density_time_derivative}
\end{equation}

Next, we split the hamiltonian into kinetic and potential parts:

\begin{equation}
    \hat{H} = \hat{T}_\mathrm{e} + \hat{V}_\mathrm{ext} + \hat{V}_\mathrm{ee}
\label{splitted_hamiltonian}
\end{equation}

One shows that both $\hat{V}_\mathrm{ext}$ and $\hat{V}_\mathrm{ee}$ commute with $\hat{n}(\vec{r}\,)$:

\begin{equation}
    \left[\hat{V}_\mathrm{ext}+\hat{V}_\mathrm{ee},\,\hat n_\mathrm{H}(\vec{r},t)\right] = 0
\label{commuting_potential_density}
\end{equation}

They are diagonal in the coordinate basis, so only the kinetic term contributes:

\begin{equation}
    \frac{\partial}{\partial t}\hat{n}_\mathrm{H}
    = \frac{\mathrm{i}}{\hbar}\left[\hat{T}_\mathrm{e},\hat{n}_\mathrm{H}\right]
\label{kinetic_term_contribution}
\end{equation}

The Heisenberg-picture current operator can now be introduced:

\begin{equation}
    \hat{\vec{j}}_\mathrm{H}(\vec{r},t)
    = \frac{1}{2m_\mathrm{e}}\sum_i\left[\hat{\vec{p}}_i\,\delta(\vec{r}-\hat{\vec{r}}_i\,)+\delta(\vec{r}-\hat{\vec{r}}_i\,)\,\hat{\vec{p}}_i\right]_\mathrm{H}
\label{heisenberg-picture_current_operator}
\end{equation}

here, $\hat{\vec{p}}_i$ is the momentum operator of particle $i$. A direct commutator calculation then yields

\begin{equation}
    \frac{\mathrm{i}}{\hbar}\left[\hat{T}_\mathrm{e},\hat n_\mathrm{H}(\vec{r},t)\right]
    = -\nabla\cdot\hat{\vec{j}}_\mathrm{H}(\vec{r},t)
\label{direct_commutator}
\end{equation}

Putting these results into equation~\eqref{density_time_derivative} gives the operator-level continuity equation:

\begin{equation}
\frac{\partial}{\partial t}\,\hat{n}_\mathrm{H}(\vec{r},t)+\nabla\cdot\hat{\vec j}_\mathrm{H}(\vec{r},t)=0
\label{continuity_operator}
\end{equation}

Taking expectation values in any Heisenberg-picture state $\ket{\Psi_\mathrm{H}}$ then yields the familiar continuity equation for the physical density $n(\vec{r},t)=\bra{\Psi_\mathrm{H}}\hat{n}_\mathrm{H}(\vec{r},t)\ket{\Psi_\mathrm{H}}$
and current $\vec{j}(\vec{r},t)=\bra{\Psi_\mathrm{H}}\hat{\vec{j}}_\mathrm{H}(\vec{r},t)\ket{\Psi_\mathrm{H}}$:

\begin{equation}
    \frac{\partial}{\partial t}\,n(\vec{r},t)+\nabla\cdot\vec{j}(\vec{r},t)=0
\label{continuity_expectation}
\end{equation}

This exact conservation law provides a fundamental constraint on the time evolution of the density and current. It also serves as one of the basic structural ingredients of time-dependent density functional theory, where the many-body dynamics are reformulated in terms of the time-dependent density $n(\vec{r},t)$.

\newpage
\section{Toward the many-body problem}
\label{sec:toward_many_body_problem}

% opening paragraph
In fact, a many-body wave function explicitly depends on all particle coordinates, making it practically impossible to calculate and store exactly due to the exponential increase in complexity with particle number.

Nevertheless, what we ultimately aim to describe are measurable properties.
These include observables like position, momentum, energy, or spin.
We can extract them from the wave function through expectation values:

\begin{equation}
    \braket{\hat{\Omega}} = \bra{\Psi}\hat{\Omega}\ket{\Psi}
\label{obvious}
\end{equation}

\subsection{An introduction to grand canonical ensemble}

Although chapter~\ref{cha:fundamentals_a} mainly focuses on the zero-temperature ground-state formalism,
it is useful to briefly record the grand-canonical framework because it provides the natural extension to finite temperature and varying particle number~\cite{gould2026ensemblization}.

Before proceeding further, it is essential to introduce the grand canonical ensemble, as it naturally accommodates quantum many-body systems exchanging particles and energy with external reservoirs.
This treatment is particularly crucial in density functional theory calculations for realistic materials, where fluctuations in particle number and thermal effects are common.

In the grand canonical ensemble, the expectation value of any observable operator $\hat{\Omega}$ is computed as follows:

\begin{equation}
    \braket{\hat{\Omega}} 
    = \frac{1}{\mathcal{Z}}\sum_{\alpha}\mathrm{e}^{-\beta(E_{\alpha}-\mu N_{\alpha})}\bra{\Psi_{\alpha}}\hat{\Omega}\ket{\Psi_{\alpha}}
\label{grand_canonical}
\end{equation}

here, the summation runs over all quantum states labeled by $\alpha$,
and $\Psi_{\alpha}$ represents the many-body quantum states of the system, characterized by their energies $E_{\alpha}$ and particle numbers $N_{\alpha}$.
The inverse temperature $\beta$ is defined as $\beta = 1/(k_B T)$, with $T$ being the temperature and $k_B$ the Boltzmann constant.
The chemical potential $\mu$ controls particle exchange between the system and its reservoir.
Therefore, the partition function $\mathcal{Z}$ serves as a normalization factor and is defined as:

\begin{equation}
    \mathcal{Z} = \sum_{\alpha} \mathrm{e}^{-\beta(E_{\alpha}-\mu N_{\alpha})}
\label{partition_function}
\end{equation}

In quantum many-body theory and calculation, employing the grand canonical ensemble allows systematic inclusion of particle number fluctuations and thermal equilibrium effects, providing a robust and realistic statistical-mechanical description essential for accurate modeling of complex quantum materials.

In what follows, we mainly focus on the ground-state setting: zero temperature and fixed particle number; while the grand-canonical form is noted for completeness and later extensions.

\subsection{From wave function to practical functionals}

Let us stay at zero temperature and fixed particle number $N$.
Consider the operator for the measurement $\sum_{i,j,\cdots}\hat\Omega\left(x_i,x_j,\cdots\right)$,
and the ground-state $N$-body wave function $\Psi(x_1,x_2,\cdots,x_N)$.
By inserting resolutions of the identity in the coordinate basis for all particles, $I=\int \mathrm{d}x\,\ket{x}\bra{x}$~\cite{shankar2017quantum},
we can represent the expectation value as:

\begin{equation}
\begin{aligned}
    \braket{\hat\Omega} 
    &= \bra{\Psi}\hat\Omega\ket{\Psi} \\
    &= \iint\cdots\int \mathrm{d}x_1\mathrm{d}x_2\cdots\mathrm{d}x_N\,\Psi^*(x_1,x_2,\cdots,x_N)\sum_{i,j,\cdots}\hat\Omega\left(x_i,x_j,\cdots\right)\Psi(x_1,x_2,\cdots,x_N)
\end{aligned}
\label{observables_N-body}
\end{equation}

Therefore, we may regard the expectation value of the observable as a functional of the many-body wave function $\Psi$, and write

\begin{equation}
\Omega[\Psi] = \bra{\Psi}\hat{\Omega}\ket{\Psi}
\label{observables_functional}
\end{equation}

here, the square brackets $[\Psi]$ (rather than parentheses) emphasize that $\Omega$ is a functional; it is a mapping from the entire wave function $\Psi$ to a scalar.

A functional simply assigns a number to a function. In our case of equation~\eqref{observables_N-body}, its form is a known integral. However, evaluating it requires the exact many-body wave function, which we generally cannot compute or store. Moreover, integration acts as an averaging operation that smooths out microscopic fluctuations, so only the coarse-grained features of the integrand contribute to the result.

Consequently, we can discard irrelevant details of the many-body wave function and work with simplified approximations tailored to specific observables.
In practice, we often use approximate wave functions: methods such as Hartree--Fock, configuration interaction, and quantum Monte Carlo follow this strategy by retaining the dominant components of the wave function while approximating or sampling the rest.

A more radical question is whether we need to work with explicit wave functions at all, or whether the expectation value in equation~\eqref{observables_N-body} can be reformulated entirely in terms of other, more tractable quantities.

To address this question, it is natural to start from the energy, since the ground-state is determined variationally through the hamiltonian expectation value.

In the case of this approximation, for a given hamiltonian $\bra{\Psi}\hat{H}\ket{\Psi}$, we can regard it as depending on the form of the kinetic energy operator $\bra{\Psi}\hat{T}\ket{\Psi}$, external potential operator $\bra{\Psi}\hat{V}_\mathrm{ext}\ket{\Psi}$, and the interaction operator $\bra{\Psi}\hat{V}_\mathrm{int}\ket{\Psi}$.

To specify the quantum dynamics and interactions of our system within the Born--Oppenheimer approximation,
we work at the fixed nuclear coordinates $\{\vec{R}_I\}$ and write the clamped nuclei hamiltonian as:

\begin{equation}
    \hat{H}_\mathrm{BO}(\{\vec{R}_I\})
    = -\frac{1}{2}\sum_i\nabla_i^2
    + \sum_i v_{\mathrm{ext}}(\hat{\vec{r}}_i\,)
    + \sum_{i<j}\frac{1}{|\vec{r}_i - \vec{r}_j\,|}
    + E_\mathrm{NN}(\{\vec{R}_I\})
\label{BO_hamiltonian}
\end{equation}

In the present electronic many-body problem within the Born--Oppenheimer approximation, 
the nuclear coordinates are clamped, and the nuclear-nuclear repulsion becomes a negligible constant. 
We can identify the interaction operator as the electron-electron Coulomb repulsion~\eqref{electrons_Coulomb},
and the external potential acting on the electrons is the attractive electron-nucleus interaction, which can be written as a sum of one-body terms~\eqref{Electron-nucleus_one-body_form}:

\begin{equation}
\begin{aligned}
    \hat{V}_\mathrm{int} &\coloneqq \hat{V}_{\mathrm{ee}}
    = \sum_{i<j}\frac{1}{|\vec{r}_i-\vec{r}_j\,|} \\
    \hat{V}_\mathrm{ext} &\coloneqq \hat{V}_{\mathrm{eN}}
    =\sum_i v_\mathrm{ext}(\hat{\vec{r}}_i\,)
\end{aligned}
\label{reduced_operator_definition}
\end{equation}

Using the expression for the electron kinetic energy operator in equation~\eqref{electrons_kinetic}, we arrive at:

\begin{equation}
\begin{aligned}
    \bra{\Psi}\hat{H}\ket{\Psi} &= \bra{\Psi}\hat{T}\ket{\Psi} + \bra{\Psi}\hat{V}_\mathrm{ext}\ket{\Psi} + \bra{\Psi}\hat{V}_\mathrm{int}\ket{\Psi} \\
    &= \bra{\Psi}-\frac{1}{2}\sum_i\nabla_i^2\ket{\Psi}
    + \bra{\Psi}\sum_i v_{\mathrm{ext}}(\vec{r}_i\,)\ket{\Psi}
    + \bra{\Psi}\sum_{i<j}\frac{1}{|\vec{r}_i-\vec{r}_j\,|}\ket{\Psi}
\end{aligned}
\label{approx}
\end{equation}

Having treated both the kinetic energy operator and the interaction term as fixed, specifying the external potential $v_\mathrm{ext}(\vec{r}\,)$ uniquely determines the hamiltonian $\hat{H}$.
For any normalized many-body state $\ket{\Psi}$, we define the energy expectation value as:

\begin{equation}
    E \coloneqq \bra{\Psi}\hat{H}\ket{\Psi}
\label{energy_expectation_value_definition}
\end{equation}

When $\ket{\Psi}$ is chosen to be an eigenstate of $\hat{H}$, denoted by $\ket{\Psi_n}$,
the corresponding eigenvalue $E_n$ and eigenstate are determined by the Schrödinger equation:

\begin{equation}
    \hat{H}\ket{\Psi_n}=E_n\ket{\Psi_n}
\label{schrodinger_eigenvalue_equation}
\end{equation}

Accordingly, one may view the spectrum and the eigenstates as depending on the external potential through the hamiltonian:

\begin{equation}
    v_\mathrm{ext}(\vec{r}\,) \longmapsto \left\{\ket{\Psi_n\left[v_\mathrm{ext}(\vec{r}\,)\right]}, E_n\left[v_\mathrm{ext}(\vec{r}\,)\right]\right\}_{n}
\label{spectrum_mapping_from_external_potential}
\end{equation}

here, $E_n\left[v_\mathrm{ext}\right]$ denotes the energy eigenvalue associated with the stationary state $\ket{\Psi_n\left[v_\mathrm{ext}\right]}$, namely:

\begin{equation}
    E_n\left[v_\mathrm{ext}(\vec{r}\,)\right] \coloneqq
    \bra{\Psi_n\left[v_\mathrm{ext}(\vec{r}\,)\right]}\hat{H}\ket{\Psi_n\left[v_\mathrm{ext}(\vec{r}\,)\right]}
\label{eigenenergy_as_expectation_value}
\end{equation}

The grand-canonical ensemble admits a compact operator formulation in terms of a trace over the many-body Hilbert space.
Invoking the cyclic invariance of the trace, we obtain:

\begin{equation}
    \braket{\hat{\Omega}}
    = \frac{1}{\mathcal{Z}}\,\operatorname{Tr}\,\left[\mathrm{e}^{-\beta\left(\hat{H}-\mu \hat{N}\right)}\,\hat{\Omega}\right]
\label{grand_canonical_expectation_trace}
\end{equation}

with the grand-canonical partition function:

\begin{equation}
    \mathcal{Z}
    =\sum_{\alpha}\mathrm{e}^{-\beta\left(E_{\alpha}-\mu N_{\alpha}\right)} 
    =\operatorname{Tr}\,\left[\mathrm{e}^{-\beta\left(\hat{H}-\mu \hat{N}\right)}\right]
\label{grand_partition_function_trace}
\end{equation}

Note that the sum over states in equation~\eqref{grand_canonical} is nothing but a trace over the many-body Hilbert space.  By using the cyclic invariance of the trace, we can therefore recast both the expectation value and the partition function in the compact operator form:

\begin{equation}
    \braket{\hat\Omega}
    = \frac{1}{\mathcal{Z}}\sum_\alpha \mathrm{e}^{-\beta(E_\alpha-\mu N_\alpha)}\bra{\Psi_\alpha}\hat\Omega\ket{\Psi_\alpha} 
\label{grand_canonical_trace}
\end{equation}

Equivalently, by expanding the trace in the eigenbasis of $\hat{H}$, equation~\eqref{grand_canonical_expectation_trace} reduces to the sum representation in equation~\eqref{grand_canonical_trace}.

In principle, all physical observables derive from the many-body Schrödinger equation~\eqref{tdse_manybody_coord}, yet the exponential growth of the Hilbert-space dimension renders direct diagonalization intractable for realistic systems.

If the kinetic energy operator and Coulomb interaction are fixed, the ground-state total energy becomes a functional of the external potential we introduced in equation~\eqref{spectrum_mapping_from_external_potential}.

A central challenge is that evaluating the grand-canonical expectation in equation~\eqref{grand_canonical_trace} remains highly nontrivial, as the trace involves high-dimensional integrals and noncommuting operators.

Our aim is to express the expectation values—which are, in principle, highly complex functionals—in terms of a single simple function, so that practical approximations become possible.

Next, recall the general expression for an observable in a many-body system: equation~\eqref{observables_N-body}. If the observable is a local one-body operator, we have:

\begin{equation}
\begin{aligned}
    \braket{\hat\Omega_1}
    &= \bra{\Psi}\hat\Omega_1\ket{\Psi} \\
    &= \iint\cdots\int \mathrm{d}x_1\mathrm{d}x_2\cdots\mathrm{d}x_N\,\Psi^*(x_1,x_2,\cdots,x_N)\sum_{i}\hat\Omega_1(x_i)\Psi(x_1,x_2,\cdots,x_N) \\
    &= \int\mathrm{d}x_1\Omega_1(x_1)N\int\mathrm{d}x_2\cdots\mathrm{d}x_N\,\Psi^*(x_1,x_2,\cdots,x_N)\Psi(x_1,x_2,\cdots,x_N)
\end{aligned}
\label{observables_one-body}
\end{equation}

We can find that this integral naturally factors into two components: the one-body operator kernel $\Omega_1(x_1)$,
and the corresponding one-body density~\cite{engel2011density,mazziotti2012two}:

\begin{equation}
    n(x_1)=N\int\mathrm{d}x_2\cdots\mathrm{d}x_N\Psi^*(x_1,x_2,\cdots,x_N)\Psi(x_1,x_2,\cdots,x_N)
\label{one-body_density}
\end{equation}

Accordingly, the ground-state expectation value of the one-body operator can be written as:

\begin{equation}
\langle\hat{\Omega}_1\rangle=\int\mathrm{d}x_1\,\Omega_1(x_1)\,n(x_1)
\label{expectation_one_body_operator}
\end{equation}

here, $x$ denotes the one-particle coordinate.
If $x$ refers to the spatial coordinate only, we write $x=\vec{r}$ and denote the usual spatial density as $n(\vec{r}\,)$.
More generally, in many-body notation $x$ denotes a combined space-spin coordinate, and we write $x=(\vec{r},\sigma)$.
In this case, $n(x)=n(\vec{r},\sigma)$ is the spin-resolved one-body density.

In many condensed matter physics applications, we work with the spin-summed density.
To recover the spin-summed spatial density, we write a discrete sum over the spin index:

\begin{equation}
    n(\vec{r}\,)=\sum_{\sigma} n(\vec{r},\sigma)
\label{sum_spatial_density}
\end{equation}

or use an integral notation for the spin index,

\begin{equation}
    n(\vec{r}\,)=\int\mathrm{d}\sigma\,n(\vec{r},\sigma)
\label{integral_spatial_density}
\end{equation}

Both notations express the same idea: the spatial density is obtained by summing over the spin degree of freedom.

%%

Throughout this quantum-mechanical discussion, for simplicity, when we write $\int\mathrm{d}x$, it means $\sum_{\sigma}\int\mathrm{d}\vec{r}$.
For spin-independent external potentials, we write $v_\mathrm{ext}(x)=v_\mathrm{ext}(\vec{r}\,)$.

The density is also constrained by particle number conservation.
For a spin-resolved density $n(x)$, we define the particle number functional as:

\begin{equation}
    N[n]\coloneqq\int\mathrm{d}x\,n(x)
\label{particle_number_functional}
\end{equation}

and in the ground-state setting with fixed particle number we require:

\begin{equation}
    N[n]=N
\label{particle_number_constraint_density}
\end{equation}

% Equivalently, for the spin-summed spatial density, the same condition reads
% $\int \mathrm{d}\vec{r}\,n(\vec{r}\,)=N$.
Equivalently, for the spin-summed spatial density, the same condition takes the form $\int \mathrm{d}\vec{r}\,n(\vec{r}\,)=N$.
This constraint will be enforced explicitly when we later vary an energy functional with respect to $n$.

We also note that not every non-negative function is an admissible density.
In general, a candidate $n(x)$ must be compatible with an antisymmetric $N$-body wave function and the chosen external potential.
In what follows, we keep these representability requirements implicit,
and we work within the set of admissible densities, and equation~\eqref{particle_number_constraint_density} provides the particle number constraint.

%%

The expectation value reduces to:

\begin{equation}
    \braket{\hat\Omega_1}=\int\mathrm{d}x_1\,\Omega_1(x_1)n(x_1)
\label{observables_one-body_density}
\end{equation}

Thus, rather than computing the full many-body integral, we need only determine the one-body density $n(x_1)$.

Next, consider the expectation value of the electron-electron Coulomb interaction energy:

\begin{equation}
\begin{aligned}
    \braket{\hat{V}_\mathrm{ee}}
    &= \iint\cdots\int \mathrm{d}x_1\mathrm{d}x_2\cdots\mathrm{d}x_N\,\Psi^*(x_1,x_2,\cdots,x_N)\sum_{i,j>i}\frac{1}{|\vec{r}_i-\vec{r}_j\,|}\Psi(x_1,x_2,\cdots,x_N) \\
    &= \int\mathrm{d}x_1\mathrm{d}x_2\,\frac{1}{2|\vec{r}_1-\vec{r}_2\,|}N(N-1)\int\mathrm{d}x_3\cdots\mathrm{d}x_N\,\Psi^*(x_1,x_2,\cdots,x_N)\Psi(x_1,x_2,\cdots,x_N)
\end{aligned}
\label{coulomb_interaction_energy}
\end{equation}

here, the two-body density naturally appears:

\begin{equation}
    n_2(x_1,x_2)=N(N-1)\int\mathrm{d}x_3\cdots\mathrm{d}x_N\,\Psi^*(x_1,x_2,\cdots,x_N)\Psi(x_1,x_2,\cdots,x_N)
\label{two-body_density}
\end{equation}

So that the expectation value of the Coulomb interaction energy reduces to:

\begin{equation}
    \braket{\hat{V}_\mathrm{ee}}=\frac{1}{2}\int\mathrm{d}x_1\mathrm{d}x_2\,
    \frac{n_2(x_1,x_2)}{|\vec{r}_1-\vec{r}_2\,|}
\label{observables_two-body_density}
\end{equation}

Similarly, the expectation value of the kinetic energy takes the form:

\begin{equation}
\begin{aligned}
    T
    &= -\iint\cdots\int \mathrm{d}x_1\mathrm{d}x_2\cdots\mathrm{d}x_N\,
    \Psi^*(x_1,x_2,\cdots,x_N)\sum_i\frac{\nabla_i^2}{2}\Psi(x_1,x_2,\cdots,x_N) \\
    &= -N\iint\cdots\int \mathrm{d}x_1\mathrm{d}x_2\cdots\mathrm{d}x_N\,
    \Psi^*(x_1,x_2,\cdots,x_N)\frac{\nabla_1^2}{2}\Psi(x_1,x_2,\cdots,x_N) \\
    &= -\frac{1}{2}\int\dif{x_1}\,
    \left[\nabla_{x'_1}^{2}\,\rho(x_1,x'_1)\right]_{x'_1=x_1}
\end{aligned}
\label{kinetic_energy}
\end{equation}

We can also separate this integral via the matrix of density:

\begin{equation}
    \rho(x_1,x'_1)
    = N\int\cdots\int\mathrm{d}x_2\cdots\mathrm{d}x_N\,\Psi^*(x_1,x_2,\cdots,x_N)\Psi(x'_1,x_2,\cdots,x_N) 
\label{density_matrix}
\end{equation}

and get the simplified expectation value of the kinetic energy operator:

\begin{equation}
    T = -\frac{1}{2}\int\dif{x_1}\,\left[\nabla_{x'_1}^{2}\,\rho(x_1,x'_1)\right]_{x'_1=x_1}
\label{kinetic_energy_density_matrix}
\end{equation}

We begin by insisting on a single functional variable $Q$ such that every observable of interest $\hat\Omega_i$ can be written in the form $\hat\Omega_i[Q]$.
Having identified $Q$, we then derive the explicit mappings $\hat\Omega_i[Q]$ for each target observable.
Finally, we develop a practical algorithm or an approximation scheme to compute $Q$ itself, thereby rendering all observables $\hat\Omega_i[Q]$ tractable.

Given a time-dependent many-body wave function $\Psi(x_1,x_2,\cdots,x_N; t)$. 
We denote the corresponding ground-state ($g\coloneqq 0$) many-body wave function as 
$\Psi_g(x_1,x_2,\cdots,x_N; t)$.
We can also denote the excited states by an index $\alpha$ ($\alpha\geq1$), 
as $\Psi_{\alpha}(x_1,x_2,\cdots,x_N; t)$.

Then, we can describe the ground-state one-body density:

\begin{equation}
    n_g(x_1,t)=N\iint\cdots\int\dif{x_2}\,\dif{x_3}\cdots\dif{x_N}\,\Psi_g^*(x_1,x_2,\cdots,x_N;t)\,\Psi_g(x_1,x_2,\cdots,x_N;t) % \rho_g(x_1,t)
\label{ground_one-body_density}
\end{equation}

When this system makes transitions to excited states, the one-body transition density is obtained by integrating other coordinates:

\begin{equation}
    f_n^\mathrm{eh}(x_1,t)=\iint\cdots\int\dif{x_2}\,\dif{x_3}\cdots\dif{x_N}\,\Psi_g^*(x_1,x_2,\cdots,x_N;t)\,\Psi_n(x_1,x_2,\cdots,x_N;t)
\label{transition_one-body}
\end{equation}

Equivalently, the reduced integral reads: 

\begin{equation}
f_n^\mathrm{eh}(x_1,t)=\int\Psi_g^*(\vec{x};t)\,\Psi_n(\vec{x};t)\prod^N_{j=2}\dif{x_j}
\label{one_body_transition_quantity_compact}
\end{equation}

This marginal integral represents the one-body transition density from the ground-state $g$ to the $n$-th excited state.

In the presence of an oscillating electromagnetic field, we introduce the Bohr frequency:

\begin{equation}
\omega_{ng}=(E_n-E_g)/\hbar
\label{Bohr_frequency}
\end{equation}

The corresponding phase factor associated with the excitation energy is $\mathrm{e}^{-\mathrm{i}\omega_{ng}t}$.
We retain $\hbar$ explicitly in these phase relations.
For a time-independent hamiltonian, we separate the time dependence of each stationary state as:

\begin{equation}
    \ket{\Psi_j(t)}=\mathrm{e}^{-\mathrm{i}E_jt/\hbar}\ket{\Phi_j},
    \qquad j\in\{g,n\}
\label{time_dep_eigenstate}
\end{equation}

here, $\ket{\Phi_j}$ is the time-independent many-body eigenstate with energy $E_j$.
Substituting these states into equation~\eqref{one_body_transition_quantity_compact},
we obtain:

\begin{equation}
    f_n^\mathrm{eh}(x_1,t)
    =\mathrm{e}^{-\mathrm{i}\omega_{ng}t}f_n^\mathrm{eh}(x_1,0)
\label{mod_transition_one-body}
\end{equation}

here, the phase-independent part follows from the coordinate projections of the stationary states:

\begin{equation}
    f_n^\mathrm{eh}(x_1,0)
    =\int\braket{\Phi_g|\vec{x}\,}\braket{\vec{x}\,|\Phi_n}
    \prod_{j=2}^{N}\dif{x_j}
\label{transition_amplitude}
\end{equation}

The time dependence is therefore contained in a single phase factor,
while the transition amplitude is determined by the stationary states.

Subsequently, we recall a one-body operator $\hat{\Omega}^{(1)}$ with kernel $\Omega_1(x_1)$,
the one-body density $n(x_1)$,
the one-body density matrix $\rho(x_1,x'_1)$,
and the pair density $n_2(x_1,x_2)$.
Then, the spatial parts are contained in $x_1,x_2$:

\begin{equation}
\begin{aligned}
    \braket{\hat{\Omega}^{(1)}} &= \int \mathrm{d}x_1\,\Omega_1(x_1)\,n(x_1)\\
    T &= -\frac{1}{2}\int \mathrm{d}x_1\,
    \left[\nabla_{x'_1}^{2}\,\rho(x_1,x'_1)\right]_{x'_1=x_1}\\
    \braket{\hat{V}_\mathrm{ee}} &= \frac{1}{2}\int \mathrm{d}x_1\,\mathrm{d}x_2\,
    \frac{n_2(x_1,x_2)}{|\vec{r}_1-\vec{r}_2\,|}
\end{aligned}
\label{one-body_operator_with_kernel}
\end{equation}

For the present discussion, we adopt $\hat{\Omega}$ to denote a typical energy-like operator composed of one-body, kinetic, and interaction contributions.
Its expectation value then follows as:

\begin{equation}
\begin{aligned}
    \braket{\hat{\Omega}} % \braket{\Omega_1}
    &= \int \mathrm{d}x_1\, \Omega_1(x_1)\, n(x_1) \\
    &\quad - \frac{1}{2}\int \mathrm{d}x_1\,
    \left[\nabla_{x'_1}^{2}\,\rho(x_1,x'_1)\right]_{x'_1=x_1} \\
    &\quad + \frac{1}{2}\int \mathrm{d}x_1\,\mathrm{d}x_2\,
    \frac{n_2(x_1,x_2)}{|\vec{r}_1-\vec{r}_2\,|}
\end{aligned}
\label{mod_obs_1}
\end{equation}

% Compact Q -> lead to density functional

Up to now, we do not have closed-form expressions for the objects appearing in equation~\eqref{mod_obs_1}, and evaluating them still requires many-body information beyond practical reach.
This motivates the introduction of a compact descriptor: rather than working with the full wave function $\Psi(x_1,\cdots,x_N)$, we seek a simpler variable $Q$ that is nevertheless sufficient to determine the observables of interest.

More precisely, let $\{\hat{\Omega}_i\}_{i=1,2,\cdots}$ denote the operators representing the target physical quantities, and let $\Omega_i[\Psi]$ denote their ground-state expectation values regarded as functionals of the wave function.
Our goal is to find a single descriptor $Q$ such that, once $Q$ is known, all these observables can be evaluated without explicit reference to the wave function $\Psi$:

\begin{equation}
    \Psi \mapsto Q,\quad\Omega_i[\Psi] = \Omega_i[Q],\quad(i=1,2,\cdots)
\label{descriptor_Q_mapping}
\end{equation}

In other words, $Q$ is chosen to retain exactly the information needed for the set $\{\Omega_i\}$, while discarding irrelevant microscopic details.

In particular, we set $Q \coloneqq n(x)$. In the following, we take the particle density $n(x)$ as the basic variable, as it provides a minimal yet powerful descriptor of the ground-state. In this statement, $x$ denotes a combined space-spin coordinate, $x=(\vec{r},\sigma)$, so that we can represent the one-body (spin-resolved) density as $n(x)$:

Firstly, the particle density $n(x)$ yields exactly all local one-body observables. For an operator of the form $\hat{\Omega}^{(1)}=\sum_{j=1}^N \Omega_1(x_j)$, one has:

\begin{equation}
    \braket{\hat{\Omega}^{(1)}}=\int \mathrm{d}x\,\Omega_1(x)\,n(x)
\label{one-body_observables}
\end{equation}

Secondly, the particle density $n(x)$ is extremely compact compared to the wave function $\Psi$.
Moreover, in the case of a non-degenerate ground-state, the Hohenberg--Kohn theorem implies a bijective correspondence between the ground-state density and the external potential.
As a result, $n(x)$ determines $v_\mathrm{ext}$ up to an additive constant.
This fixes the hamiltonian and thereby fixes the ground-state up to an overall phase.
Therefore, all ground-state observables may be regarded as density functionals, and we seek $\Omega_i[n]$ rather than $\Omega_i[\Psi]$.

Therefore, we explore the variational route for $n$. Define the energy-like functional:

\begin{equation}
    E[n;v_\mathrm{ext}]=T[n]+E_\mathrm{int}[n]+\int\mathrm{d}x\,v_\mathrm{ext}(x)n(x)
\label{energy-like_functional}
\end{equation}

and obtain the ground-state density from
\begin{equation}
    \frac{\delta}{\delta n}\left(E[n;v_\mathrm{ext}] - \mu\,N[n]\right)=0 \ \Rightarrow\ n_g(x)
\label{ground_state_density}
\end{equation}

Once $n_g$ is known, all desired observables are evaluated as $\Omega_i[n_g]$; 
for non-local one-body or two-body quantities,
this amounts to adopting controlled approximations for kinetic energy functional $T[n]$ and interaction energy functional $E_\mathrm{int}[n]$ such as the Kohn--Sham framework, which we will introduce later.

\newpage
\section{Orbital-free density functionals}
\label{sec:appendix_thomas_fermi}

\subsection{The Thomas--Fermi approach}

Now, we introduce the Thomas--Fermi approach as a first and purely density-based model for the kinetic energy.

The basic idea is to treat the electrons locally as a uniform, spin-degenerate Fermi gas.
At each point $\vec{r}$, the local density $n(\vec{r}\,)$ fixes a local Fermi momentum $p_F(\vec{r}\,)$, which allows us to estimate the kinetic energy density from the occupied Fermi sphere in momentum space.

For a spin-degenerate system, the phase-space counting gives:

\begin{equation}
    n(\vec{r}\,) = \frac{2}{h^3}\frac{4\pi}{3}\,p_F^3(\vec{r}\,)
\label{TF_density_pf_intro}
\end{equation}

and the local Fermi momentum is related to the local single-particle energy through:

\begin{equation}
    \mu = \frac{p_F^2(\vec{r}\,)}{2m}+v(\vec{r}\,)
\label{TF_mu_local_intro}
\end{equation}

Combining these two expressions yields a local relation between $n(\vec{r}\,)$ and $v(\vec{r}\,)$.
This relation also leads to an explicit approximation for the kinetic energy functional in terms of $n(\vec{r}\,)$ alone, which we introduce next.

To close the Thomas--Fermi picture, we treat $v(\vec{r}\,)$ as an electrostatic potential generated self-consistently by the electron density.
In atomic units, it satisfies the Poisson equation:

\begin{equation}
    \nabla^2 v(\vec{r}\,) = -4\pi\,n(\vec{r}\,)
\label{TF_poisson_equation}
\end{equation}

Eliminating $p_F(\vec{r}\,)$ between \eqref{TF_density_pf_intro} and \eqref{TF_mu_local_intro}, we obtain:

\begin{equation}
    n(\vec{r}\,) = \frac{1}{3\pi^2}\,[2(\mu - v(\vec{r}\,))]^{3/2}
\label{TF_density_muV}
\end{equation}

Substituting this expression into equation~\eqref{TF_poisson_equation},
We arrive at the Thomas--Fermi equation:

\begin{equation}
    \nabla^2 v(\vec{r}\,)=-\frac{4}{3\pi}\,[2(\mu-v(\vec{r}\,))]^{3/2}
\label{thomas_fermi_equation}
\end{equation}

This nonlinear equation determines the potential and the density self-consistently within the Thomas--Fermi approximation.

To connect the Thomas--Fermi construction with density functional theory, we rewrite the total energy in a density functional form, denoted by $E_{\mathrm{TF}}[n;v_\mathrm{ext}]$. In this sense, the Thomas--Fermi model provides an explicit approximation to the universal functional.

From the Thomas--Fermi approximation, we express the total energy as a density functional composed of a statistical kinetic term $T_\mathrm{TF}[n]$, an interaction term $E_\mathrm{int}[n]$, and an external contribution $E_\mathrm{ext}[n]$ describing the coupling to the external potential $v_\mathrm{ext}(\vec{r}\,)$.

Furthermore, there is also a nucleus-nucleus term $E_\mathrm{NN}$, but it does not depend on the electron density. Therefore, we do not consider this term at the moment.

Accordingly, we write the Thomas--Fermi total energy functional:

\begin{equation}
    E_{\mathrm{TF}}[n;v_\mathrm{ext}] = T_\mathrm{TF}[n] + E_\mathrm{int}[n] + E_\mathrm{ext}[n]
\label{TF_total_energy_functional}
\end{equation}

According to the local form of the external potential energy functional~\eqref{functional_Eext_local_form}, and the external contribution is given by the coupling between $n(\vec{r}\,)$ and $v_\mathrm{ext}(\vec{r}\,)$.
For fixed nuclear positions, the external potential $v_\mathrm{ext}(\vec{r}\,)$ is a prescribed function of $\vec{r}$ defined in equation~\eqref{one-body_Electron-nucleus}. Therefore, the external contribution $E_\mathrm{ext}[n]$ is linear in the electron density $n(\vec{r}\,)$.

We next approximate the interaction energy functional by its classical Coulomb form associated with the density distribution.
Using the electron-electron Coulomb repulsion~\eqref{electrons_Coulomb}, we introduce the Hartree term:

\begin{equation}
\begin{aligned}
    E_\mathrm{H}[n] &\coloneqq E_\mathrm{int}[n]-E_{\Delta\mathrm{int}}[n] \\
    &= \frac{1}{2}\iint\dif{\vec{r}}\,\dif{\vec{r}\,'}\,\frac{n(\vec{r}\,)n(\vec{r}\,')}{|\vec{r}-\vec{r}\,'|}
\label{hartree_term_functional}
\end{aligned}
\end{equation}

here, $E_{\Delta\mathrm{int}}[n]$ collects the nonclassical part of the interaction energy beyond the Hartree term.
In the Hartree approximation, we set $E_{\Delta\mathrm{int}}[n]=0$.

Therefore, we introduce a Hartree approximation at the level of the interaction energy by replacing $E_\mathrm{int}[n]$ with its classical counterpart $E_\mathrm{H}[n]$:

\begin{equation}
    E_\mathrm{H}^{\mathrm{int}}[n]
    \coloneqq T[n]
    + \int \dif^3{r}\,v_\mathrm{ext}(\vec{r}\,)n(\vec{r}\,)
    + E_\mathrm{H}[n]
\label{hartree_interaction_level_energy}
\end{equation}

here, $T[n]$ is the interacting kinetic energy functional evaluated on the minimizing many-body state $\Psi[n]$ that yields the density $n$.
In other words, this replacement affects only the interaction energy, while the kinetic energy functional remains $T[n]$ at this stage.

We next consider the kinetic energy functional.
For a homogeneous spin-unpolarized electron gas with density $n$,
the Fermi wave vector satisfies~\cite{simon2013oxford}:

\begin{equation}
    k_\mathrm{F} = \left(3\pi^2 n\right)^{\frac{1}{3}}
\label{Fermi_wave_vector_homogeneous_gas}
\end{equation}

and the corresponding kinetic energy density reads:

\begin{equation}
    t(n) = \frac{3}{10}\,k_\mathrm{F}^2\,n
    = \frac{3}{10}\left(3\pi^2\right)^{\frac{2}{3}}n^{\frac{5}{3}}
\label{kinetic_energy_density_homogeneous_gas}
\end{equation}

Within the local density approximation,
we apply this homogeneous result locally by replacing $n$ with $n(\vec{r}\,)$ and integrating over space.
Accordingly, this yields the Thomas--Fermi kinetic energy functional~\cite{thomas1927calculation,fermi1927metodo,mi2023orbital}:

\begin{equation}
    T_\mathrm{TF}[n] = C_1\int\dif{\vec{r}}\,n(\vec{r}\,)^{\frac{5}{3}}
\label{functional_TTF}
\end{equation}

with the Thomas--Fermi constant:

\begin{equation}
    C_1 \coloneqq \frac{3}{10}\left(3\pi^2\right)^{\frac{2}{3}} 
\label{Thomas_Fermi_constant_C1}
\end{equation}

This expression holds for a spin-unpolarized system.

Therefore, we arrive at the full Thomas--Fermi energy functional as a functional of the electron density~\cite{garcia2007efficient,xu2024recent}.
The nucleus-nucleus contribution $E_\mathrm{NN}$ is omitted since it does not depend on $n(\vec{r}\,)$.

\begin{equation}
\begin{aligned}
    E_\mathrm{TF}[n]
    &= C_1\int\dif{\vec{r}}\,n(\vec{r}\,)^{\frac{5}{3}}
    + \frac{1}{2}\iint\dif{\vec{r}}\,\dif{\vec{r}\,'}\,\frac{n(\vec{r}\,)n(\vec{r}\,')}{|\vec{r}-\vec{r}\,'|}
    + \int\dif{\vec{r}}\,v_\mathrm{ext}(\vec{r}\,)\,n(\vec{r}\,) \\
    &= C_1\int\dif{\vec{r}}\,n(\vec{r}\,)^{\frac{5}{3}}
    + E_\mathrm{H}[n] + \int\dif{\vec{r}}\,v_\mathrm{ext}(\vec{r}\,)\,n(\vec{r}\,)
\end{aligned}
\label{full_energy_functional}
\end{equation}

This expression gives the total energy in the Thomas--Fermi approximation as a functional of the electron density $n(\vec{r}\,)$ for a given external potential $v_\mathrm{ext}(\vec{r}\,)$.

The ground-state of the system is characterized by the density that minimizes the total energy under the constraint of a fixed particle number.
This is the variational principle for the ground-state energy in density functional theory.

At finite temperature, the equilibrium density is obtained by minimizing the Helmholtz free energy functional~\cite{mermin1965thermal}.
In the present zero temperature setting, the relevant quantity is the total energy functional $E_\mathrm{TF}[n]$.

To determine the ground-state density, we minimize the Thomas--Fermi energy functional under the particle number constraint:

\begin{equation}
    \int\dif{\vec{r}}\,n(\vec{r}\,)=N
\label{constraint_particle_number}
\end{equation}

Using the method of Lagrange multipliers, we require the first variation of the constrained functional to vanish.

The Lagrange multiplier $\lambda$ enforces the particle number constraint and can be identified with the chemical potential:

\begin{equation}
\begin{aligned}
    0 &= \delta\left\{E_\mathrm{TF}[n]-\lambda\left(\int\dif{\vec{r}}\,n(\vec{r}\,)-N\right)\right\}\\
    &= \delta\left\{T_\mathrm{TF}[n]+E_\mathrm{int}[n]+E_\mathrm{ext}[n]-\lambda\left(\int\dif{\vec{r}}\,n(\vec{r}\,)-N\right)\right\}
\end{aligned}
\label{total_energy_variation}
\end{equation}

Here $\lambda$ enforces the particle number constraint and plays the role of the chemical potential.

We define the Hartree potential as a functional of density:

\begin{equation}
    v_\mathrm{H}([n];\vec{r}\,)\coloneqq\int\dif{\vec{r}\,'}\,\frac{n(\vec{r}\,')}{|\vec{r}-\vec{r}\,'|}
\label{Hartree_potential_definition}
\end{equation}

The corresponding first variations of the three contributions can be evaluated explicitly.

For the Thomas--Fermi kinetic term, we obtain:

\begin{equation}
    \delta T_\mathrm{TF}[n]
    = \int\dif{\vec{r}}\,\frac{5}{3}C_1\,n(\vec{r}\,)^{\frac{2}{3}}\,\delta n(\vec{r}\,)
\label{variation_TTF}
\end{equation}

For the interaction term in the Hartree approximation, the variation introduces the Hartree potential and reads:

\begin{equation}
    \delta E_\mathrm{int}[n]
    = \int\dif{\vec{r}}\,v_\mathrm{H}([n];\vec{r}\,)\,\delta n(\vec{r}\,)
\label{variation_Eint}
\end{equation}

For the external contribution, the variation follows directly from the linear coupling to the external potential:

\begin{equation}
    \delta E_\mathrm{ext}[n]
    = \int\dif{\vec{r}}\,v_\mathrm{ext}(\vec{r}\,)\,\delta n(\vec{r}\,)
\label{variation_Eext}
\end{equation}

Substituting equations~\eqref{variation_TTF}--\eqref{variation_Eext} into equation~\eqref{total_energy_variation}, we obtain:

\begin{equation}
    \int\dif{\vec{r}}\,\left\{\frac{5}{3}C_1\,n(\vec{r}\,)^{\frac{2}{3}}+v_\mathrm{H}([n];\vec{r}\,)+v_\mathrm{ext}(\vec{r}\,)-\lambda\right\}\delta n(\vec{r}\,)=0
\label{total_energy_minimization}
\end{equation}

Since the density variation $\delta n(\vec{r}\,)$ is arbitrary, the integrand must vanish pointwise, which yields the Thomas--Fermi equation:

\begin{equation}
    \frac{5}{3}C_1\,n(\vec{r}\,)^{\frac{2}{3}}+v_\mathrm{H}([n];\vec{r}\,)+v_\mathrm{ext}(\vec{r}\,)=\lambda
\label{Thomas_Fermi_equation}
\end{equation}

This is convenient to introduce the effective potential:

\begin{equation}
    v(\vec{r}\,)\coloneqq v_\mathrm{ext}(\vec{r}\,)+v_\mathrm{H}([n];\vec{r}\,)
\label{effective_potential_definition}
\end{equation}

Equation~\eqref{Thomas_Fermi_equation} then becomes:

\begin{equation}
    \frac{5}{3}C_1\,n(\vec{r}\,)^{\frac{2}{3}}+v(\vec{r}\,)=\lambda
\label{Thomas_Fermi_equation_compact}
\end{equation}

Solving for the density, we obtain:

\begin{equation}
    n(\vec{r}\,)=\frac{1}{3\pi^2}\left[2\left(\lambda-v(\vec{r}\,)\right)\right]^{\frac{3}{2}}
\label{density_TF_from_potential}
\end{equation}

The above expression applies in the region where $\lambda>v(\vec{r}\,)$, and one takes $n(\vec{r}\,)=0$ otherwise.

To derive a differential form, we use the Poisson equation.
For the nuclear external potential, we have:

\begin{equation}
    \nabla^2 v_\mathrm{ext}(\vec{r}\,)=4\pi\sum_I Z_I\,\delta\left(\vec{r}-\vec{R}_I\,\right)
\label{Poisson_external_potential}
\end{equation}

For the Hartree potential, we have:

\begin{equation}
    \nabla^2 v_\mathrm{H}([n];\vec{r}\,)=-4\pi n(\vec{r}\,)
\label{Poisson_Hartree_potential}
\end{equation}

Combining equations~\eqref{effective_potential_definition}, \eqref{Poisson_external_potential}, and \eqref{Poisson_Hartree_potential}, we find:

\begin{equation}
    \nabla^2 v(\vec{r}\,)=4\pi\sum_I Z_I\,\delta\left(\vec{r}-\vec{R}_I\,\right)-4\pi n(\vec{r}\,)
\label{Poisson_effective_potential}
\end{equation}

Substituting equation~\eqref{density_TF_from_potential} into equation~\eqref{Poisson_effective_potential}, we arrive at the differential form of the Thomas--Fermi equation:

\begin{equation}
    \nabla^2 v(\vec{r}\,)=4\pi\sum_I Z_I\,\delta\left(\vec{r}-\vec{R}_I\,\right)-\frac{4}{3\pi}\left[2\left(\lambda-v(\vec{r}\,)\right)\right]^{\frac{3}{2}}
\label{differential_form_Thomas_Fermi_equation}
\end{equation}

This procedure constitutes the Thomas--Fermi method.

\subsection{Limitations of the Thomas--Fermi approach}

For a many-body system,
it is instructive to contrast the exact density functional with the Thomas--Fermi approximation.
In density functional theory, the total energy can be written as a functional of the electron density~\cite{hohenberg1964inhomogeneous,jones2015density}.

Let us express the exact total energy as a functional of the density:

\begin{equation}
    E[n]=T[n]+E_\mathrm{H}[n]+E_\mathrm{xc}^\mathrm{int}[n]+\int\dif{\vec{r}}\,n(\vec{r}\,)v_\mathrm{ext}(\vec{r}\,)
\label{functional_exact_total_energy}
\end{equation}

Here $T[n]$ denotes the exact kinetic energy functional and $E_\mathrm{H}[n]$ is the Hartree energy functional.
The exchange-correlation functional $E_\mathrm{xc}^\mathrm{int}[n]$ collects all nonclassical interaction effects beyond the Hartree term.
It includes the exchange contribution originating from the antisymmetry of the many-electron wave function and the correlation contribution accounting for the remaining many-body correlations.

In view of the definition of the universal Hohenberg--Kohn functional in equation~\eqref{functional_HK_definition},
it is convenient to introduce the exchange-correlation part of the interaction energy functional via:

\begin{equation}
E_\mathrm{xc}^{\mathrm{int}}[n]\coloneqq E_\mathrm{int}[n]-E_\mathrm{H}[n]
\label{Exc_int_definition}
\end{equation}

Equivalently, the universal Hohenberg--Kohn functional can be decomposed as:

\begin{equation}
\begin{aligned}
    E_\mathrm{HK}[n] &= T[n] + E_\mathrm{int}[n] \\
    &= T[n]+E_\mathrm{H}[n]+E_\mathrm{xc}^\mathrm{int}[n]
\end{aligned}
\label{HK_decomposition_Hartree_xc}
\end{equation}

Accordingly, the exact total energy functional can be written in the following form:

\begin{equation}
    E[n]=E_\mathrm{HK}[n]+\int\dif{\vec{r}}\,n(\vec{r}\,)\,v_\mathrm{ext}(\vec{r}\,)
\label{functional_exact_total_energy_HK_form}
\end{equation}

Recall the full Thomas--Fermi energy functional in equation~\eqref{full_energy_functional}.
Let us compare with the exact decomposition in equation~\eqref{functional_exact_total_energy},
the Thomas--Fermi functional neglects the exchange correlation part of the interaction energy functional $E_\mathrm{xc}^\mathrm{int}[n]$.
This omission can lead to sizable errors for total energies even in simple atoms~\cite{schwinger1980thomas, schwinger1981thomas}.

A first improvement is to approximate $E_\mathrm{xc}^\mathrm{int}[n]$ by a local exchange functional derived from the homogeneous electron gas.
This approximation is commonly associated with the Dirac exchange functional~\cite{dirac1930note}:

\begin{equation}
    E_\mathrm{xc}^\mathrm{int}[n]\approx E_\mathrm{x}^\mathrm{D}[n]
\label{Dirac_exchange_approximation}
\end{equation}

Since the Dirac exchange functional is given by:

\begin{equation}
    E_\mathrm{x}^\mathrm{D}[n]\coloneqq -C_2\int\dif{\vec{r}}\,n(\vec{r}\,)^{\frac{4}{3}}
\label{Dirac_exchange_functional}
\end{equation}

here, the corresponding constant reads:

\begin{equation}
    C_2 \coloneqq \frac{3}{4}\left(\frac{3}{\pi}\right)^{\frac{1}{3}}
\label{Dirac_exchange_constant}
\end{equation}

At this level, correlation effects are not included in $E_\mathrm{x}^\mathrm{D}[n]$.

Including the Dirac exchange term, we obtain the Thomas--Fermi Dirac functional:

\begin{equation}
    E_\mathrm{TFD}[n]
    \coloneqq C_1\int\dif{\vec{r}}\,n(\vec{r}\,)^{\frac{5}{3}}
    + E_\mathrm{H}[n]
    - C_2\int\dif{\vec{r}}\,n(\vec{r}\,)^{\frac{4}{3}}
    + \int\dif{\vec{r}}\,n(\vec{r}\,)\,v_\mathrm{ext}(\vec{r}\,)
\label{TFD_functional}
\end{equation}

The Thomas--Fermi Dirac approximation improves total energies by incorporating an explicit exchange contribution.
However, the kinetic energy is still treated at the Thomas--Fermi level, which is often insufficient for molecules and solids.

A systematic gradient correction to the kinetic energy can be introduced through the von Weizsäcker form~\cite{yang1986gradient}.
The corresponding von Weizsäcker kinetic energy functional is defined as:

\begin{equation}
    T_\mathrm{vW}[n]\coloneqq \frac{1}{2}\int\dif{\vec{r}}\,\left|\nabla\sqrt{n(\vec{r}\,)}\right|^2
\label{von_Weizsaecker_kinetic}
\end{equation}

By adding this gradient correction, we arrive at the Thomas--Fermi Dirac von Weizsäcker model:

\begin{equation}
    E_\mathrm{TFDW}[n]
    \coloneqq C_1\int\dif{\vec{r}}\,n(\vec{r}\,)^{\frac{5}{3}}
    + T_\mathrm{vW}[n]
    + E_\mathrm{H}[n]
    - C_2\int\dif{\vec{r}}\,n(\vec{r}\,)^{\frac{4}{3}}
    + \int\dif{\vec{r}}\,n(\vec{r}\,)\,v_\mathrm{ext}(\vec{r}\,)
\label{TFDW_functional}
\end{equation}

These extensions provide early orbital free approximations beyond the Thomas--Fermi approach~\cite{lieb1981thomas}.
They also clarify that improving the kinetic energy functional and the exchange correlation functional is essential for a quantitatively reliable description of realistic interacting electron systems.

\newpage
\section{Toward electromagnetic response in matter}
\label{sec:toward_electromagnetic_response_in_matter}

% opening paragraph
Chapter~\ref{cha:fundamentals_a} focused on the ground-state electronic problem.
There, the central question was how electrons arrange themselves under a given external potential.
However, optical properties do not follow from the ground-state alone.
When an electromagnetic field is applied,
it drives the electrons away from their static equilibrium and triggers a series of responses.

Therefore, before introducing the dielectric function, we need a language that describes both the applied field and the feedback of matter to that field.
That language is electrodynamics.
The present section does not aim to review a complete lecture of electrodynamics;
instead, we only extract the ingredients required for the discussion of the dielectric function, the Kramers--Kronig relations, and the optical properties in chapter~\ref{cha:fundamentals_b}~\cite{landau2013electrodynamics,griffiths2023introduction}.
We begin with electromagnetic fields in matter, and then gradually move toward a frequency-dependent response function that can be connected to the electronic structure.

We use SI units for the macroscopic electrodynamic equations.
The dielectric function is dimensionless and denotes the relative dielectric response.
We use the angular frequency $\omega$ with the time dependence $\exp(-\mathrm{i}\,\omega t)$,
and the photon energy is $\hbar\omega$.

\subsection{Maxwell equations in matter}

Let us start by the universal structure of the problem.
When an electromagnetic field enters a medium, it drives charge redistribution and induces currents inside the medium, while these induced sources further modify the field itself.
Therefore, the optical response of matter should be described within the framework of electrodynamics~\cite{landau2013electrodynamics,griffiths2023introduction}.

At this stage, we do not explore how a specific medium responds;
instead, we first identify the general field equations that govern electromagnetic dynamics in matter.
They provide the universal backbone of the problem, while the constitutive relations later encode the medium-specific information.

Firstly, we write the macroscopic Maxwell equations for the electric field $\vec{E}(\vec{r},t)$ and magnetic field $\vec{B}(\vec{r},t)$ as:

\begin{equation}
\left\{\begin{aligned}
    \nabla\cdot\vec{E}(\vec{r},t) &= \frac{\rho(\vec{r},t)}{\varepsilon_0} 
    &&\space \text{Gauss's law for electric field} \\
    \nabla\cdot\vec{B}(\vec{r},t) &= 0 
    &&\space \text{Gauss's law for magnetic field} \\
    \nabla\times\vec{E}(\vec{r},t) &= -\frac{\partial}{\partial t}\vec{B}(\vec{r},t) 
    &&\space \text{Faraday's law} \\
    \nabla\times\vec{B}(\vec{r},t) &= \mu_0\vec{J}(\vec{r},t)+\mu_0\varepsilon_0\frac{\partial}{\partial t}\vec{E}(\vec{r},t)
    &&\space \text{Ampère--Maxwell law}
\end{aligned}\right.
\label{maxwell_equations_with_total_sources}
\end{equation}

here, $\rho(\vec{r},t)$ and $\vec{J}(\vec{r},t)$ denote the total charge density and total current density, respectively.

These equations already imply a local conservation law.
More specifically, we take the divergence of the Ampère--Maxwell law,
and use Gauss's law for the electric field to obtain the continuity equation for charge at the macroscopic level~\cite{landau2013electrodynamics}:

\begin{equation}
    \frac{\partial}{\partial t}\rho(\vec{r},t)+\nabla\cdot\vec{J}(\vec{r},t)=0
\label{macroscopic_charge_continuity_equation}
\end{equation}

here, the conservation principle discussed in section~\ref{sec:time_evolution_quantum_mechanics} reappears in electrodynamic language.

Secondly, for media, it proves useful to decompose the sources into free and bound contributions:

\begin{equation}
\begin{aligned}
    \rho(\vec{r},t) &= \rho_\mathrm{free}(\vec{r},t)+\rho_\mathrm{bound}(\vec{r},t) \\
    \vec{J}(\vec{r},t) &= \vec{J}_\mathrm{free}(\vec{r},t)+\vec{J}_\mathrm{bound}(\vec{r},t)
\end{aligned}
\label{free_and_bound_source_decomposition}
\end{equation}

here, the free part describes externally supplied or mobile macroscopic sources,
while the bound part describes the response generated internally by the medium.

For discussing the bound electric response, we introduce the polarization field $\vec{P}(\vec{r},t)$ as:

\begin{equation}
    \rho_\mathrm{bound}(\vec{r},t)\coloneqq -\nabla\cdot\vec{P}(\vec{r},t)
\label{bound_charge_from_polarization}
\end{equation}

Accordingly, we define the electric displacement field as:

\begin{equation}
    \vec{D}(\vec{r},t)\coloneqq \varepsilon_0\vec{E}(\vec{r},t)+\vec{P}(\vec{r},t)
\label{electric_displacement_field_definition}
\end{equation}

For the optical response considered here, we neglect magnetic-permeability effects and take the relative magnetic permeability as $\mu_{\mathrm{r}}(\omega)\approx1$.
The electric response is then the central object, and we keep the Maxwell equations in the form:

\begin{equation}
\left\{\begin{aligned}
    \nabla\cdot\vec{D}(\vec{r},t) &= \rho_\mathrm{free}(\vec{r},t) 
    &&\space \text{Gauss's law for electric displacement} \\
    \nabla\cdot\vec{B}(\vec{r},t) &= 0 
    &&\space \text{Gauss's law for magnetic field} \\
    \nabla\times\vec{E}(\vec{r},t) &= -\frac{\partial}{\partial t}\vec{B}(\vec{r},t) 
    &&\space \text{Faraday's law} \\
    \frac{1}{\mu_0}\nabla\times\vec{B}(\vec{r},t) &= \vec{J}_\mathrm{free}(\vec{r},t)+\frac{\partial}{\partial t}\vec{D}(\vec{r},t)
    &&\space \text{Amp\`ere--Maxwell law in matter}
\end{aligned}\right.
\label{maxwell_equations_in_matter}
\end{equation}

Consequently, we can now characterize the structure of the problem.
The Maxwell equations themselves are universal,
while the medium dependence enters through the induced polarization $\vec{P}(\vec{r},t)$, and therefore through the electric displacement field $\vec{D}(\vec{r},t)$.

In many optical problems, one considers the response to an external electromagnetic perturbation in the absence of free charges and free currents inside the medium.
Under this condition, equation~\eqref{maxwell_equations_in_matter} shows that the central issue is no longer the field equations themselves,
but how the medium develops polarization under the applied electric field.
Once the relation between $\vec{P}(\vec{r},t)$ and $\vec{E}(\vec{r},t)$ is specified, the electromagnetic response becomes well defined.
This is precisely the point where macroscopic electrodynamics meets microscopic electronic structure,
and it also provides the starting point for introducing the dielectric function~\cite{landau2013electrodynamics,lima2022usefulness,bernasconi2011first}.

\subsection{Polarization, electric displacement, and constitutive relations}

We have established the universal field equations in the previous subsection.
However, those equations do not explore how matter responds to the applied field.
To close the problem, we still need a relation that connects the induced polarization to the electric field.
This relation is the constitutive relation, and it is here that the medium-specific information enters~\cite{landau2013electrodynamics, lima2022usefulness}.

Here, the central object is the polarization field $\vec{P}(\vec{r},t)$.
It describes how the medium rearranges its internal charges under the applied electric field.
Once $\vec{P}(\vec{r},t)$ is known, equation~\eqref{electric_displacement_field_definition} immediately gives the electric displacement field $\vec{D}(\vec{r},t)$.
Therefore, instead of treating $\vec{D}(\vec{r},t)$ as an independent quantity, we may regard it as a convenient way to encode the electric response of matter.

Firstly, we start from the most general linear-response form.
For any time $t$, the polarization may depend not only on the electric field at the same point,
but also on fields acting earlier and in a surrounding spatial region.
Therefore, before introducing the simplified form used in the present work,
we write the constitutive relation in its general nonlocal form in space and time~\cite{royo2019first}:

\begin{equation}
    P_\alpha(\vec{r},t) =
    \varepsilon_0\sum_\beta\int\dif^3{r'}
    \int_{-\infty}^{t}\dif{t'}
    \,\chi_{\alpha\beta}(\vec{r},t;\vec{r}\,',t')\,E_\beta(\vec{r}\,',t')
\label{general_nonlocal_constitutive_relation}
\end{equation}

here, $\chi_{\alpha\beta}(\vec{r},t;\vec{r}\,',t')$ describes the contribution of the electric field acting at $(\vec{r}\,',t')$ to the polarization that develops later at $(\vec{r},t)$.
In this sense, the medium does not have to respond only at one point and one instant;
instead, it may spread the response over space and retain memory in time.
At the same time, the upper limit $t$ in equation~\eqref{general_nonlocal_constitutive_relation} keeps the causal order explicit: the perturbation acts first, and the response follows.

For the present work, we focus on the linear response in the long-wavelength regime.
Under this approximation, spatial nonlocality drops out, and the constitutive relation retains only the memory of the earlier perturbation.
Therefore, it reduces to the time-domain constitutive relation:

\begin{equation}
    P_\alpha(\vec{r},t)=
    \varepsilon_0\sum_\beta\int_{-\infty}^{t}\dif{t'}\,\chi_{\alpha\beta}(t-t')\,E_\beta(\vec{r},t')
\label{time_domain_constitutive_relation}
\end{equation}

here, $\chi_{\alpha\beta}(t-t')$ denotes the time-dependent susceptibility tensor.
This expression demonstrates that the polarization does not follow the electric field instantaneously;
instead, the medium retains a memory of the earlier perturbation.

Optical experiments usually probe the response of a medium under an oscillating electromagnetic field.
Therefore, once we enter the linear-response regime, the frequency domain provides the natural language for discussing the optical response~\cite{han2023temperature}.
Next, we transform equation~\eqref{time_domain_constitutive_relation} to the frequency domain, and obtain:

\begin{equation}
    P_\alpha(\vec{r},\omega)=\varepsilon_0\sum_\beta\chi_{\alpha\beta}(\omega)\,E_\beta(\vec{r},\omega)
\label{frequency_domain_polarization_susceptibility_relation}
\end{equation}

here, $\chi_{\alpha\beta}(\omega)$ denotes the frequency-dependent electric susceptibility tensor.
This equation already shows that the susceptibility determines the linear electric response of the medium.

Substituting equation~\eqref{frequency_domain_polarization_susceptibility_relation} into equation~\eqref{electric_displacement_field_definition},
we then express the electric displacement field in terms of the electric susceptibility tensor:

\begin{equation}
\begin{aligned}
    D_\alpha(\vec{r},\omega)
    &=\varepsilon_0 E_\alpha(\vec{r},\omega)+P_\alpha(\vec{r},\omega) \\
    &=\varepsilon_0\sum_\beta\left[\delta_{\alpha\beta}+\chi_{\alpha\beta}(\omega)\right]E_\beta(\vec{r},\omega)
\end{aligned}
\label{electric_displacement_from_susceptibility}
\end{equation}

It now shows that the combination $\delta_{\alpha\beta}+\chi_{\alpha\beta}(\omega)$ naturally enters the linear electric response.

This observation directly motivates the definition of the dielectric tensor $\varepsilon_{\alpha\beta}(\omega)$:

\begin{equation}
    \varepsilon_{\alpha\beta}(\omega)
    \coloneqq\delta_{\alpha\beta}+\chi_{\alpha\beta}(\omega)
\label{dielectric_tensor_definition}
\end{equation}

Then, the frequency-domain constitutive relation takes the compact form:

\begin{equation}
    D_\alpha(\vec{r},\omega)
    = \varepsilon_0\sum_\beta\varepsilon_{\alpha\beta}(\omega)E_\beta(\vec{r},\omega)
\label{frequency_domain_constitutive_relation_with_dielectric_tensor}
\end{equation}

The dielectric tensor $\varepsilon_{\alpha\beta}(\omega)$ thus provides a compact description of the linear electric response.
Starting from the polarization field $\vec{P}(\vec{r},t)$,
we introduced the electric susceptibility tensor $\chi_{\alpha\beta}(\omega)$,
and then arrived at the dielectric tensor $\varepsilon_{\alpha\beta}(\omega)$.
Therefore, this tensor emerges naturally from the constitutive relation that closes Maxwell electrodynamics in matter.

Besides its directional dependence, the dielectric response is also dynamical.
In linear response theory, causality constrains this dynamical character,
which means that the polarization at a given time can only depend on the electric field applied at earlier times.
This causal structure becomes important in section~\ref{sec:causality_derivation_basic},
where we derive the Kramers--Kronig relations from the analytic structure of the response.

\subsection{Anisotropy and the dielectric tensor}

The previous subsection presented the dielectric tensor $\varepsilon_{\alpha\beta}(\omega)$ as a compact description of the linear electric response.
We now examine how symmetry constrains its form~\cite{capper2019optical}.
When the medium responds identically along all directions, the dielectric tensor reduces to a scalar form.
If the response changes with direction, the full tensor form must be retained.

Firstly, for an isotropic medium,
the dielectric tensor $\varepsilon_{\alpha\beta}(\omega)$ reduces to a scalar dielectric function $\varepsilon(\omega)$~\cite{born2013principles}:

\begin{equation}
    \varepsilon_{\alpha\beta}(\omega)
    = \varepsilon(\omega)\delta_{\alpha\beta}
\label{isotropic_dielectric_tensor}
\end{equation}

here, the scalar dielectric function $\varepsilon(\omega)$ already determines the full linear electric response,
since the isotropic medium does not distinguish one Cartesian direction from another.
The tensor structure, therefore, no longer carries additional directional information.

Secondly, crystalline media usually distinguish directions.
In that case, anisotropy survives, and we must retain the full tensor form.
When the crystal symmetry allows a diagonal dielectric tensor in the chosen Cartesian axes,
the dielectric tensor $\varepsilon_{\alpha\beta}(\omega)$ takes the principal-axis form~\cite{nye1985physical}:

\begin{equation}
    \left[\varepsilon_{\alpha\beta}(\omega)\right] = \begin{pmatrix}
        \varepsilon_{xx}(\omega) & 0 & 0 \\
        0 & \varepsilon_{yy}(\omega) & 0 \\
        0 & 0 & \varepsilon_{zz}(\omega)
    \end{pmatrix}
\label{principal_axis_dielectric_tensor}
\end{equation}

here, each diagonal component describes the dielectric response along one principal axis,
and anisotropy enters through the difference among these components.

Here, this directional dependence plays a central role~\cite{slavich2024exploring}.
In chapter~\ref{cha:fundamentals_b}, we compare the optical response within the structural plane and perpendicular to that plane.
When the two in-plane directions are symmetry-equivalent,
we introduce the in-plane dielectric function $\varepsilon_{\varparallel}(\omega)$ and the out-of-plane dielectric function $\varepsilon_{\perp}(\omega)$~\cite{nye1985physical,guo2021complete,matthes2016influence}:

\begin{equation}
    \left[\varepsilon_{\alpha\beta}(\omega)\right] = \begin{pmatrix}
        \varepsilon_{\varparallel}(\omega) & 0 & 0 \\
        0 & \varepsilon_{\varparallel}(\omega) & 0 \\
        0 & 0 & \varepsilon_{\perp}(\omega)
    \end{pmatrix}
\label{inplane_outofplane_dielectric_tensor}
\end{equation}

here, the in-plane dielectric function $\varepsilon_{\varparallel}(\omega)$ describes the response to an electric field within the structural plane,
while the out-of-plane dielectric function $\varepsilon_{\perp}(\omega)$ describes the response to an electric field normal to that plane.
This notation directly prepares the comparison between the in-plane and out-of-plane dielectric spectra and optical properties~\cite{guo2021complete,slavich2024exploring}.

This anisotropy also propagates to other optical properties derived from the dielectric tensor~\cite{sujith2023large,slavich2024exploring}.
Once the dielectric response changes with direction, the refractive index, the extinction coefficient, the reflectivity, and the energy-loss spectrum also inherit that directional dependence.
Therefore, the tensor structure does not only refine the notation.
It also determines how we organize and interpret the later optical results.

At this stage, the dielectric tensor $\varepsilon_{\alpha\beta}(\omega)$ has been introduced as the quantity that encodes the directional linear electric response of the medium.
When anisotropy is present, its tensor form makes this directional dependence explicit.

Once the principal directions of the dielectric tensor $\varepsilon_{\alpha\beta}(\omega)$ are chosen as the working axes, the electric field $\vec{E}(\omega)$ can be decomposed as:

\begin{equation}
    \vec{E}(\omega)=\sum_{\lambda}E_{\lambda}(\omega)\,\hat{e}_{\lambda}
\label{electric_field_decomposition_principal_directions}
\end{equation}

where $\hat{e}_{\lambda}$ denotes the unit vector along the principal direction $\lambda$.
In the same basis, the induced polarization $\vec{P}(\omega)$ takes the form:

\begin{equation}
    \vec{P}(\omega)=\varepsilon_0\sum_{\lambda}\chi_{\lambda}(\omega)\,E_{\lambda}(\omega)\,\hat{e}_{\lambda}
\label{polarization_decomposition_principal_directions}
\end{equation}

Accordingly, we denote the dielectric response along each principal direction by $\varepsilon_{\lambda}(\omega)$ and express:

\begin{equation}
    \varepsilon_{\lambda}(\omega)=1+\chi_{\lambda}(\omega)
\label{direction_resolved_dielectric_response}
\end{equation}

This direction-resolved form provides the natural bridge from the tensor description to the dielectric function discussed in section~\ref{sec:from_electronic_response_to_dielectric_function}.

\newpage
\section{Causality and the dielectric response}
\label{sec:causality_derivation_basic}

Section~\ref{sec:causality_and_kramers_kronig_relations} uses causality to connect the real and imaginary parts of the dielectric function.
Here, we establish the analytic structure of the response and derive the full-frequency Kramers--Kronig relations.

\subsection{Causality and analytic structure of response functions}

Firstly, to express causality explicitly, 
we start from the time-domain linear relation between
the polarization field $P_{\alpha}(\vec{r},t)$ and
the electric field $E_{\beta}(\vec{r}\,',t')$~\cite{dressel2002electrodynamics}:

\begin{equation}
    P_{\alpha}(\vec{r},t) = 
    \varepsilon_0\sum_{\beta}\int_{-\infty}^{\infty}\dif{t'}
    \int\dif^3{r'}\,\chi_{\alpha\beta}(\vec{r},\vec{r}\,',t-t')
    E_{\beta}(\vec{r}\,',t')
\label{general_linear_response_polarization}
\end{equation}

here, $\chi_{\alpha\beta}(\vec{r},\vec{r}\,',t-t')$ denotes the time-domain electric susceptibility tensor.
This equation characterizes the polarization response at position $\vec{r}$ and time $t$ to an electric field applied at position $\vec{r}\,'$ and time $t'$.

Therefore, causality requires that the medium responds only to earlier external driving fields~\cite{stefanski2022analytical}:

\begin{equation}
    \chi_{\alpha\beta}(\vec{r},\vec{r}\,',t-t')=0
    \qquad \text{for $t<t'$}
\label{causal_condition_susceptibility_kernel}
\end{equation}

This condition defines a retarded response kernel~\cite{landau2013electrodynamics}.
It excludes any polarization response before the external electric field acts.

In the present work,
we consider the optical long-wavelength limit,
where the spatial nonlocality of the dielectric response becomes negligible~\cite{agranovich2013crystal}.
Afterwards, the linear response then reduces to a local frequency-dependent tensor relation~\cite{franta2023dispersion}:

\begin{equation}
    P_{\alpha}(\omega) =
    \varepsilon_0\sum_{\beta}\chi_{\alpha\beta}(\omega)
    \,E_{\beta}(\omega)
\label{frequency_domain_polarization_response}
\end{equation}

here, $\chi_{\alpha\beta}(\omega)$ denotes the frequency-domain electric susceptibility tensor.
It characterizes the response amplitude for each Cartesian component at angular frequency $\omega$.
Meanwhile, it also describes the phase delay of the response at the same frequency~\cite{koutserimpas2024time}.

We then express the dielectric tensor $\varepsilon_{\alpha\beta}(\omega)$ via the electric susceptibility tensor $\chi_{\alpha\beta}(\omega)$~\cite{franta2020symmetry}:

\begin{equation}
    \varepsilon_{\alpha\beta}(\omega) =
    \delta_{\alpha\beta}
    + \chi_{\alpha\beta}(\omega)
\label{dielectric_tensor_from_susceptibility}
\end{equation}

Subsequently, the causal condition in equation~\eqref{causal_condition_susceptibility_kernel}
gives the frequency-domain susceptibility in the form:

\begin{equation}
\chi_{\alpha\beta}(\omega) =
    \int_{0}^{\infty}
    \chi_{\alpha\beta}(t)
    \,\mathrm{e}^{\,\mathrm{i}\omega t}
    \dif{t}
\label{frequency_domain_susceptibility_from_causality}
\end{equation}

This representation reveals the analytic structure directly.
Since the time integral starts from $t=0$, the response function remains analytic in the upper half of the complex frequency plane~\cite{nussenzveig1972causality}.

For a real time-domain response kernel, the susceptibility $\chi_{\alpha\beta}(\omega)$
and dielectric tensor $\varepsilon_{\alpha\beta}(\omega)$ further satisfy the following condition:

\begin{equation}
\left\{\begin{aligned}
    \chi_{\alpha\beta}(-\omega)
    &= \chi_{\alpha\beta}^{*}(\omega) \\
    \varepsilon_{\alpha\beta}(-\omega)
    &= \varepsilon_{\alpha\beta}^{*}(\omega)
\end{aligned}\right.
\label{reality_condition_response_functions}
\end{equation}

Therefore, in the principal-axis representation,
for each complex dielectric component $\varepsilon_{\lambda\lambda}(\omega)$,
the real part $\varepsilon_{1,\lambda\lambda}(\omega)$ is an even function of frequency,
whereas the imaginary part $\varepsilon_{2,\lambda\lambda}(\omega)$ is an odd function:

\begin{equation}
\left\{\begin{aligned}
    \varepsilon_{1,\lambda\lambda}(-\omega)
    &= \varepsilon_{1,\lambda\lambda}(\omega) \\
    \varepsilon_{2,\lambda\lambda}(-\omega)
    &= -\varepsilon_{2,\lambda\lambda}(\omega)
\end{aligned}\right.
\label{parity_of_dielectric_components}
\end{equation}

At this stage, the previous discussion has established the analytic structure of the dielectric response in frequency space~\cite{nussenzveig1972causality,franta2020symmetry,leon2026momentum}.
We next derive the Kramers--Kronig relations between its real and imaginary parts.

\subsection{Kramers--Kronig relations for the dielectric function}

The previous subsection established the analytic structure of the dielectric response in frequency space.
We now derive the Kramers--Kronig relations for the dielectric function~\cite{toll1956causality}.
Firstly, the causal structure in equation~\eqref{frequency_domain_susceptibility_from_causality} and the reality condition in equation~\eqref{reality_condition_response_functions}
establish a Hilbert-transform relation between the real and imaginary parts of the dielectric response~\cite{stefanski2022analytical,koutserimpas2024time,franta2020symmetry}.
Here, we consider the regular electronic response without a zero-frequency conductivity pole,
with $\varepsilon_{\alpha\beta}(\omega)\to\delta_{\alpha\beta}$ at high frequency.
For a metal, the conductive term must be separated before applying these relations in the form below.
The dielectric function in this derivation then denotes the remaining regular response.
We represent the Kramers--Kronig relations in the frequency-domain form as:

\begin{equation}
\left\{\begin{aligned}
    \varepsilon_{1,\alpha\beta}(\omega)-\delta_{\alpha\beta} &=
    \frac{1}{\pi}\mathcal{P}
    \int_{-\infty}^{\infty}\dif{\omega'}\,
    \frac{\varepsilon_{2,\alpha\beta}(\omega')}{\omega'-\omega} \\
\varepsilon_{2,\alpha\beta}(\omega) &=
    -\frac{1}{\pi}\mathcal{P}
    \int_{-\infty}^{\infty}\dif{\omega'}\,
    \frac{\varepsilon_{1,\alpha\beta}(\omega')-\delta_{\alpha\beta}}{\omega'-\omega}
\end{aligned}\right.
\label{kk_general_dielectric_tensor}
\end{equation}

here, $\mathcal{P}$ denotes the Cauchy principal value.
It defines the integral through a symmetric limiting procedure around $\omega'=\omega$~\cite{stefanski2022analytical,koutserimpas2024time}.

This expression identifies the physical roles of the two parts of the dielectric response.
The imaginary part $\varepsilon_{2,\alpha\beta}(\omega)$ characterizes the absorptive channel, 
whereas the real part $\varepsilon_{1,\alpha\beta}(\omega)$ characterizes the dispersive channel.
Consequently, causality unifies these two channels into a single complex response function.

The symmetry relations in equation~\eqref{parity_of_dielectric_components} then reduce the full-frequency integrals to the positive-frequency form in equation~\eqref{kk_positive_frequency_dielectric_component}.
The corresponding static limit and its physical interpretation are discussed in section~\ref{sec:causality_and_kramers_kronig_relations}.
